\section{The horizon: build the means of further construction}
\label{sec:horizon}

\subsection{Scientific instruments are a particularly valuable beginning}

A new instrument can change which questions are answerable. A better experiment can change which design uncertainties remain binding. A production process can then turn that knowledge into repeatable material capability. This gives a construction a second possible consequence: it may improve the means of discovering and realizing another construction.

The relationship is causal only when the successor actually depends on the new means relative to strong alternatives. A later project that uses the same codebase is not enough. Neither is a story that every useful invention somehow contributes to everything after it. Identify the specific observation, operation, representation, resource, or arrangement the predecessor supplied; remove or replace it in a fair comparison; measure the successor consequence.

Three ambitious candidate families illustrate different mechanisms. A new \emph{measurement operation} can expose an otherwise unidentifiable property. A new \emph{conversion or manufacturing process} can turn an unusable input into a qualified component. A new \emph{productive arrangement} can make complementary facilities, demand, expertise, and funding viable together. These are research directions, not completed products of this edition.

\begin{figure}[htbp]
\centering
\begin{tikzpicture}[every node/.style={font=\small,align=center},node distance=.29in and .25in,box/.style={draw=Rule,fill=Ivory,rounded corners=3pt,text width=2.32in,minimum height=.73in,inner sep=7pt},arr/.style={-{Latex[length=2.5mm]},line width=.9pt,color=Teal}]
\node[box] (observe) {\textbf{BUILD WAYS TO KNOW}\\Observations, experiments,\\representations};
\node[box,right=of observe] (produce) {\textbf{BUILD WAYS TO MAKE}\\Qualified processes, instruments,\\manufacturing};
\node[box,below=of produce] (expand) {\textbf{BUILD WAYS TO EXPAND}\\Energy, materials, logistics,\\maintainable capacity};
\node[box,below=of observe] (learn) {\textbf{IMPROVE THE BUILDER}\\Retain mechanisms and test\\fresh construction};
\draw[arr] (observe)--(produce);\draw[arr] (produce)--(expand);\draw[arr] (expand)--(learn);\draw[arr] (learn)--(observe);
\draw[Gold,dashed,line width=.6pt] ($(observe.north west)+(-.12,.12)$) rectangle ($(expand.south east)+(.12,-.12)$);
\end{tikzpicture}
\caption{The long-horizon hypothesis. Each arrow requires its own accepted construction and supported enabling relationship. The loop does not imply a fixed sequence or guaranteed growth.}
\end{figure}

\subsection{Knowledge can travel; capacity has to be made}

A useful design can be reused without rebuilding the original reasoning. This is an important property of knowledge, connected to the economic distinction between rival resources and nonrival technological knowledge \citep{romer1990}. The design's physical realization still requires materials, energy, equipment, fabrication, qualification, logistics, and time. Its rights and applicability may also limit transfer.

Consider a prospective production batch $u_t$, a vector of proposed activities. Necessary resource conditions can be written
\begin{equation}
A_tu_t\le r_t,\qquad e_t^{\mathsf T}u_t\le E_t,\qquad
h_t^{\mathsf T}u_t\le H_t.
\label{eq:physical}
\end{equation}
Here $A_t$ records material and other capacity requirements in declared compatible units, $r_t$ their available amounts, $E_t$ the energy budget over the batch interval, and $H_t$ the admissible heat-rejection budget over that interval. The coefficient vectors record the corresponding batch requirements. Scheduling, peak power, environmental conditions, quality, losses, and maintenance add further constraints; these aggregate inequalities alone are not sufficient for operation.

Intelligence can change the coefficients by inventing a different process. It can find an alternative material, reduce wasted work, improve measurement, or remove an unnecessary dependency. It cannot simply omit a term because the result would otherwise be unaffordable. A cash balance is not an energy supply, and a bill of materials is not a commissioned facility.

The ambitious target is productive infrastructure that helps construct more productive infrastructure. Qualified sensing supports process control; control supports fabrication; fabrication supports additional instruments and energy systems. Replacement parts and repair are part of this loop. A system that expands quickly while consuming an unreplaceable critical input is not self-sustaining under that boundary.

\subsection{Beyond one facility, beyond one planet}

The horizon extends to extensive solar-powered production, autonomous construction, and scientific infrastructure beyond Earth. This is not an untouched conceptual territory. NASA's current lunar technology programme explicitly combines power, resource processing, manufacturing, construction, and environmental challenges; its descriptions include tests in simulated gravity and regolith-simulant environments \citep{nasa2026}. Those established programmes and qualified components belong among the strongest alternatives to assess. Their achievements remain theirs.

The research question is which remaining constraint prevents a particular useful system from operating, and what new construction can remove it. If the constraint is a missing reference or unreliable conversion step, a precise small experiment may be more consequential than a broad architectural demonstration. If it is logistics or financing, another algorithm may be beside the point. If the incumbent already meets the purpose, the discovery programme must look for a different residual rather than invent one.

Expansion to much larger energy and material throughput would require many separately qualified advances: long-lived power systems, resource acquisition, manufacturing chains, repair, transport, thermal management, coordination, and legitimate governance. No measured result here supplies an expansion rate, deployment date, or wealth forecast. The paper's contribution is to make the route a sequence of buildable, testable dependencies rather than a leap from intelligence to an assumed outcome.

\begin{thesisbox}
\textbf{The durable ambition is an increasing ability to create useful abilities.}\par
Scientific knowledge matters through what it lets us understand and do. Productive wealth matters through worthwhile possibilities it sustains. A successful construction programme connects them by building, checking, maintaining, and improving the means.
\end{thesisbox}

\subsection{Expansion should enlarge useful choice}

Private returns can fund research and construction, but they are not the same quantity as total benefit. Competition and diffusion can reduce one firm's margin while making a capability useful to more people. A worthwhile system must remain compatible with voluntary participation, honest risk allocation, exit, and meaningful alternatives.

Construction can also create harm, concentrate control, exhaust resources, or make a community dependent on fragile infrastructure. These are consequences to investigate and govern, not external details to ignore when counting productive capacity. The relevant objective is sustained useful possibility under legitimate authority, with affected parties and material obligations represented.

There is no universal business-size ceiling in the framework. Nor is there a guarantee that bigger is better. Real resource constraints, end-use value, maintenance, coordination, environmental effects, and alternatives determine the useful next scale.

\section{Limits, falsifiers, and the next result worth earning}
\label{sec:limits}

The finite study supplies its task language, calibration coefficients, costs, times, rights flags, and construction method. It therefore tests implementation and conditional reachability, not the autonomous discovery of those objects. Its three successor tasks vary within one authored arithmetic family. Four thousand input checks do not amount to four thousand independent scientific discoveries. The conventional control ties; the study provides no measured constructor-learning advantage.

The pretrained route meets the nucleus research contract but fails the whole-cell boundary requirement. The classical constructor searches a supplied finite recipe bank; its selected source fails validation, and all six successor arms exhaust the bank and fail boundary acceptance. The same-content control explains the apparent construction-method result: zero saved candidates and no accepted advantage. A single exposed assay comparison cannot establish general construction learning.

The formal results are elementary consequences of declared transition, value, and resource assumptions. Real completeness depends on the adequacy of the declared boundary. Open environments can introduce new failures, participants can decline or change terms, and source evidence can be wrong. A finite exact optimum is an optimum only in its supplied feasible set.

The ambitious thesis would be weakened if retained construction methods failed on independent tasks, if ordinary equal-content routes explained their entire apparent advantage, if maintenance consumed the savings, or if useful outcomes could not support a workable arrangement. These are productive possible conclusions. An enterprise can still adopt a useful conventional tool without attributing to itself an improved invention process.

The next result worth earning is one complete consequential construction with independent acceptance, followed by a fresh construction whose improvement is causally attributable to what was retained. An economical accepted delivery would establish a further bridge. Repeated improvement of the constructor across independently initialized campaigns would establish a stronger one. These are concrete tests of the vision.

\section{Conclusion: the possibility after the possibility}

Intelligence becomes materially consequential when it helps make worthwhile outcomes attainable. That often requires more than solving a problem inside an inherited setup. It requires discovering the missing means, constructing a coherent whole, and learning which parts of the experience make new construction better.

The complete means may be a measurement, a representation, a process, a material capability, a resource, or a voluntary arrangement. The best construction can combine them, replace them, or make some unnecessary. What survives should improve a later ability, with its applicability and obligations visible.

The programme is ambitious because the object of improvement is the means of making progress itself. It is testable because each step names a useful outcome, a construction, a complete alternative, and evidence that the change mattered.

\begin{thesisbox}
{\Large\bfseries\color{Plum}Change the game.}\par
{\Large\bfseries\color{Plum}Build what the new game requires.}\par
{\Large\bfseries\color{Plum}Learn to build better.}\par\medskip
When we build something new, we change what can be built next.\par
\textbf{The task is to make that change worthwhile---and learn how to do it again.}
\end{thesisbox}
